\documentclass[runningheads]{llncs}
\usepackage{amsmath,amssymb}
\usepackage{graphicx}
\usepackage{booktabs}
\usepackage{multirow}
\usepackage{hyperref}
\hypersetup{hidelinks}
\usepackage{microtype}
\emergencystretch=3em

\title{A Systematic Evaluation of Prompt Design for LLM-Based Sexism
Detection}
\titlerunning{Prompt Design for LLM-Based Sexism Detection}
% Reverted to this paper's original working-placeholder title on
% 2026-08-26 at the user's explicit request, undoing the 2026-08-25
% rename to "The Limits of Prompt Engineering for LLM-Based Sexism
% Detection."

\author{Erfan Hoseingholizadeh \and Shakiba Mirbagheri \and Mina Farmanbar}
\authorrunning{Hoseingholizadeh et al.}
\institute{}

\begin{document}

\maketitle

\begin{abstract}
Prompt design offers an appealing route to sexism detection with large
language models: no labeled data, no fine-tuning, and a wide space of
role, context, and reasoning instructions to choose from. This paper
tests that route systematically, evaluating a curated grid of 18 prompt
variants across six open-weight, instruction-tuned models spanning five
model families and a 3.8--9B parameter range, on the full official EDOS
sexism-detection test split ($n=4{,}000$). Prompt design produces large,
statistically real effects: 92 of 101 variant-versus-baseline
comparisons reach significance under McNemar's test and 88 of 101 under
a paired F1-diff bootstrap, and the single largest of these tested
effects comes from prompt design alone.
But that entire space sits, for five of the six models, below what a
classical TF-IDF and logistic regression classifier already achieves on
the same data at a fraction of the computational cost; only one model
edges past it on F1, at 4--16$\times$ the inference cost of every other
prompted model, and a separate confidence-based AUC comparison crowns a
different model instead, not the F1 leader. Reading raw model output rather
than trusting aggregate scores further revealed that one model's
chain-of-thought failures split cleanly into two distinct causes,
explicit safety refusal or reasoning-token-budget exhaustion, depending
on prompt composition, and that prompt-component order matters far less
than prior work on few-shot demonstration order would predict. Prompt
engineering is a real, measurable lever on this task; it is not, on this
evidence, a substitute for adapting a model to the task's own labeled
data.

\keywords{sexism detection \and prompt engineering \and large language
models \and in-context learning \and chain-of-thought prompting \and
hate speech detection}
\end{abstract}

\section{Introduction}
\label{sec:introduction}

Online sexism is widespread, and automated detection is one of the tools
platforms and researchers use to make identifying it tractable at scale
\cite{kirk2023semeval}. The Explainable Detection of Online Sexism
(EDOS) benchmark, which this paper evaluates against throughout, was
built specifically to move that detection effort past a coarse
sexist/not-sexist label toward a fine-grained taxonomy of how sexism
shows up in text. Most computational work on EDOS, and on hate speech
detection more broadly, still relies on fine-tuning a model on labeled
examples. Instruction-tuned large language models offer a different
route: Brown et al.~\cite{brown2020gpt3} showed that a sufficiently
capable model can perform a new task from a natural-language description
and a handful of demonstrations, with no weight update at all, which
makes prompting an attractive option whenever labeled data, compute, or
both are scarce.

That option comes with its own open question: which prompt does the
job. A prompt for sexism detection can assign the model a role, define
the linguistic markers that make language sexist, supply surrounding
context about the platform and content type, ask the model to reason
before answering, or show worked examples, each a real design choice
with no established default for this specific task. Related work on
prompt engineering generally, and on hate speech detection with LLMs
specifically (Sec.~\ref{sec:related}), suggests these choices can matter
a great deal, but the evidence for sexism detection in particular has so
far come from single models evaluated on a narrow slice of that design
space, without a non-LLM baseline or formal significance testing to say
whether an apparent improvement is real.

This paper closes that gap directly. We evaluate a curated grid of 18
prompt variants, combining five independently toggleable components
(role, aspects, context, chain-of-thought, and few-shot demonstrations)
in isolation and in combination, across six open-weight, instruction-tuned
models spanning five model families and a 3.8--9B parameter range, on the
full official EDOS test split. Every variant's effect against a
zero-component baseline is tested for statistical significance with two
independent tests, and every model's best training-free variant is
compared directly against a QLoRA-tuned version of the same model, a
fine-tuned encoder, and a classical TF-IDF and logistic regression
baseline trained on the same data. We further break results down by
sexism subtype and test three robustness axes prior work has not applied
to this task: the order in which prompt components appear, the number of
few-shot demonstrations shown, and which specific demonstrations are
drawn.

The results argue for a narrower and more specific claim than either
``prompt design matters'' or ``LLMs solve sexism detection'' taken
alone. Prompt design does produce large, statistically real effects
within training-free prompting: the large majority of this study's
variant comparisons reach significance under each of two independent
tests, and the single largest effect among this study's
significance-tested comparisons comes from prompt design alone. But
that entire space sits, for five of the six evaluated models, below
what a linear classifier trained on the same labeled data already
achieves, at a fraction of the computational cost; only one model
clears that bar on F1, at a steep compute premium, and a separate
confidence-based AUC comparison crowns a different model instead.
Along the way, reading raw model output rather than trusting aggregate
scores surfaced a chain-of-thought-triggered failure mode in one model
that splits cleanly into two distinct causes depending on prompt
composition, and showed that prompt-component order, unlike few-shot
demonstration order in prior work, has a comparatively narrow effect on
these models' behavior. Sec.~\ref{sec:limitations} and
Sec.~\ref{sec:conclusion} address how far these findings generalize
beyond the setting tested here.


\section{Related Work}
\label{sec:related}

\subsection{Sexism and Hate Speech Detection Datasets}
\label{sec:related:datasets}

The EDOS dataset \cite{kirk2023semeval} evaluated throughout this paper was
built for SemEval-2023 Task 10, which split sexism detection into three
subtasks of increasing granularity: a binary sexist/not-sexist label, a
four-category classification of sexist content, and an eleven-vector
fine-grained taxonomy. Most published EDOS work addresses these subtasks by
fine-tuning a pretrained encoder on the training split, the same strategy
this paper's own encoder baseline
(Sec.~\ref{sec:methodology:baselines}) follows: shared-task submissions
ranged from BERT-based classifiers \cite{rifat2023acsmkrhr} to conventional
machine learning pipelines built on hand-engineered features
\cite{panwar2023edos}, the latter close in spirit to the TF-IDF baseline
used here. Where an LLM appears at all in this literature, it is typically
one comparison point among several fine-tuned models, not the object of a
systematic prompt-design study.

Sexism-specific datasets remain comparatively rare outside English. SWSR
\cite{jiang2021swsr} is, to our knowledge, the only comparably scaled
non-English sexism dataset built with its own dedicated lexicon, drawn from
Sina Weibo rather than the Reddit and Gab sources behind EDOS. Its
existence is a reminder that the platform and language choices behind EDOS
are not incidental: what counts as sexist language, and how directly it
gets stated, varies with both. This study inherits EDOS's English,
Reddit/Gab-sourced scope rather than testing across languages or
platforms, a scope limitation the paper's closing sections return to.

\subsection{Prompt Engineering and In-Context Learning}
\label{sec:related:prompting}

Brown et al. \cite{brown2020gpt3} showed that a sufficiently large frozen
language model can pick up a new task from a short natural-language
description and a handful of demonstrations, with no gradient update at
all. That result opened a design space, mapped by Liu et al.'s survey
\cite{liu2021pretrain}, of choices a practitioner makes when turning a
task into a prompt: how to phrase the instruction, whether to supply
demonstrations and how many, whether to ask the model to reason before
answering. The five toggleable components studied here, role, aspects,
context, chain-of-thought, and few-shot
(Sec.~\ref{sec:methodology:prompts}), occupy one corner of that space,
chosen for a single task rather than surveyed in the abstract.

Chain-of-thought prompting \cite{wei2022chain} is the best-known instance
of the reasoning component in that space, proposed for arithmetic and
multi-step logical problems rather than classification. Sexism detection
is not obviously the kind of task chain-of-thought was built for, and the
results here bear that out unevenly: only two of six models
(Sec.~\ref{sec:results:grid}) show a large, consistent benefit from it,
which argues against assuming a chain-of-thought instruction helps by
default on a new task rather than measuring it directly.

Order matters too. Lu et al. \cite{lu2022fantastically} found that
permuting the order of few-shot demonstrations alone can move a model from
near-state-of-the-art to near-random accuracy on standard classification
benchmarks, with no change to which demonstrations are shown. The order
sweep in this paper (Sec.~\ref{sec:methodology:sweeps}) asks a related but
distinct question: not the order of demonstrations within a fixed prompt,
but the order of entire prompt components (context, chain-of-thought,
aspects, few-shot) relative to each other. The effect is far more
contained than Lu et al.'s general finding: reordering shifts F1 for only
one of six models (Sec.~\ref{sec:results:sensitivity}), suggesting
component-order sensitivity is not a universal property of these models
the way demonstration-order sensitivity appears to be.

\subsection{Large Language Models for Hate Speech and Content Moderation}
\label{sec:related:llm-moderation}

The finding closest in spirit to this paper's own baseline comparison
(Sec.~\ref{sec:results:baselines}) comes from Plaza-del-Arco et al.
\cite{plazadelarco2023label}, who ran several LLMs zero-shot across a
range of text classification tasks and reported that none of them rivaled
even simple supervised baselines. This paper reaches a similar conclusion
by a different route: an 18-variant grid of systematically combined
prompt components per model, rather than zero-shot prompting alone, so
the gap to a linear classifier persists after substantial prompt
engineering, not only in its absence. A comparable pattern shows up in
hate speech detection specifically. Bernal-Beltr\'an et al.
\cite{bernalbeltran2025spanish} compared fine-tuned Spanish encoders
against zero- and few-shot LLM prompting on a hate speech corpus and found
the fine-tuned encoders consistently ahead, with few-shot prompting
narrowing but not closing the gap. Neither study applies the statistical
significance testing or the multi-model breadth used here
(Sec.~\ref{sec:methodology:stats}), but both point the same direction:
for classification tasks like these, a general-purpose LLM's prompting
still lags a model trained on the task's own labels.

Safety alignment introduces a failure mode distinct from accuracy.
R\"ottger et al.'s XSTest \cite{rottger2023xstest} is a purpose-built,
adversarial suite of safe-sounding prompts designed to trigger unwarranted
refusals, and it shows the phenomenon exists in several production models.
The llama3.1 refusal behavior documented in this paper
(Sec.~\ref{sec:results:llama-refusal}) was not sought out through an
adversarial test suite; it surfaced as a byproduct of ordinary sexism
classification once role and aspect framing were dropped from the prompt,
on genuinely sexist and non-sexist text rather than deliberately
borderline phrasing. It is a naturally occurring instance of the same
over-refusal phenomenon XSTest was built to probe, arising from an
ordinary classification pipeline rather than a red-teaming exercise.

\subsection{Fine-Tuned Baselines: Encoders and Parameter-Efficient
Adaptation}
\label{sec:related:baselines}

The encoder baseline in this study
(Sec.~\ref{sec:methodology:baselines}) fine-tunes RoBERTa
\cite{liu2019roberta}, whose central finding, that BERT's original
training recipe left real performance on the table, shows that a
fine-tuning baseline's strength depends on training choices as much as
architecture. The QLoRA baseline builds on two more recent ideas: LoRA
\cite{hu2021lora}, which adapts a frozen model by training a small pair of
low-rank matrices injected into its attention layers instead of updating
all its weights, and QLoRA \cite{dettmers2023qlora}, which quantizes the
frozen base weights to 4 bits so the same adaptation fits on a single GPU
even for a multi-billion-parameter model. Together they make it practical
to fine-tune all six of this paper's prompted models rather than only the
smallest, which is what makes the QLoRA-versus-prompting comparison in
Sec.~\ref{sec:results:baselines} possible at this scale.


\section{Methodology}
\label{sec:methodology}

\subsection{Prompt-Design Space}
\label{sec:methodology:prompts}

We study five independently toggleable prompt components layered on a
common base instruction (classify the input text as containing sexism or
not): a system-level \textbf{role} specification identifying the model as
a sexism-annotation assistant; a bulleted list of \textbf{aspects}
(seven linguistic markers associated with sexist language: objectification,
gender stereotyping, hostile language, benevolent sexism, dehumanization,
double standards, and dismissiveness); a short \textbf{context} block
naming the platform, content type, and language; \textbf{chain-of-thought}
(T) instructions asking the model to reason before answering; and
\textbf{few-shot} (F) demonstrations.

Rather than the full $2^5{=}32$-cell factorial, we evaluate a curated grid
of 18 variants (V1--V18): the bare baseline with every component off (V1),
each component in isolation (V2--V6), and twelve further combinations
spanning pairwise through five-way interactions, culminating in all five
components active at once (V18). Table~\ref{tab:main-grid}
(Sec.~\ref{sec:results:grid}) lists the exact component composition of
every variant alongside its results, so it is not repeated here.

Active (non-role) components are concatenated into the prompt's user turn
in a fixed order (context, chain-of-thought, aspects, few-shot) ahead
of the input text; the role component, when active, is instead injected as
a separate system-role chat message. Non-chain-of-thought variants append
an explicit trailing \texttt{ANSWER:} cue immediately before generation,
forcing an immediate two-token \texttt{YES}/\texttt{NO} decision;
chain-of-thought variants omit this cue and instead instruct the model to
reason first and conclude with its own \texttt{ANSWER: YES} /
\texttt{ANSWER: NO} line, so the forced cue never pre-empts the model's own
reasoning (Sec.~\ref{sec:experimental:generation} details the corresponding
generation-budget difference between the two). Few-shot variants sample
$k{=}2$ demonstrations per class (2 sexist, 2 not-sexist; 4 total),
seeded ($\text{seed}{=}42$) and drawn without replacement from the EDOS
training split (Sec.~\ref{sec:experimental:dataset}), formatted as
\texttt{TEXT}/\texttt{ANSWER} pairs.

\subsection{Model Roster}
\label{sec:methodology:models}

We evaluate six open-weight, instruction-tuned LLMs spanning five model
families and a 3.8--9B parameter range (Table~\ref{tab:models}). Every
model is loaded from a Hugging Face Hub revision pinned to a fixed commit
hash, so a later upstream update to a repository's default branch cannot
silently change which weights a rerun evaluates. Two models
(llama3.1, gemma2) are gated, requiring an accepted Hugging Face license;
the remaining four are openly accessible. qwen3 is the only reasoning-tuned
model in the roster: its chat template exposes a runtime
\texttt{enable\_thinking} toggle, wired to each variant's own
chain-of-thought flag, giving it genuine extended reasoning on
chain-of-thought variants rather than the brief prose reasoning the other
five models produce from instructions alone.

\begin{table}[t]
\centering
\footnotesize
\caption{Evaluated model roster. All Hugging Face Hub revisions are
pinned to a fixed commit (Sec.~\ref{sec:methodology:models}).}
\label{tab:models}
\begin{tabular}{lllc}
\toprule
Model & Hugging Face ID & Params & Reasoning \\
\midrule
mistral   & mistralai/Mistral-7B-Instruct-v0.3 & 7B   & \\
phi3      & microsoft/Phi-3-mini-4k-instruct   & 3.8B & \\
llama3.1  & meta-llama/Llama-3.1-8B-Instruct   & 8B   & \\
qwen2.5   & Qwen/Qwen2.5-7B-Instruct           & 7B   & \\
gemma2    & google/gemma-2-9b-it               & 9B   & \\
qwen3     & Qwen/Qwen3-8B                      & 8B   & \checkmark \\
\bottomrule
\end{tabular}
\end{table}

\subsection{Fine-Tuned and Classical Baselines}
\label{sec:methodology:baselines}

All three baselines below train on the EDOS training split and are scored
on the same held-out test split as the prompting grid
(Sec.~\ref{sec:experimental:dataset}), for direct comparability with the
prompting results (Sec.~\ref{sec:results:baselines}).

\textbf{QLoRA.} For each of the six prompted models we additionally
fine-tune a QLoRA adapter: 4-bit NF4-quantized frozen base weights (bf16
compute dtype) with LoRA adapters ($r{=}16$, $\alpha{=}32$, dropout
$0.05$) targeting the attention projection matrices
(\texttt{q\_proj}, \texttt{k\_proj}, \texttt{v\_proj}, \texttt{o\_proj}).
Training runs 6 epochs (learning rate $2{\times}10^{-4}$, batch size 16,
max sequence length 512) on the training split, oversampled to a balanced
50/50 class ratio before tokenization: the SFT objective (next-token
prediction on a single \texttt{" YES"}/\texttt{" NO"} completion) has no
built-in class weighting, and, trained directly on EDOS's natural
$\sim$76/24 imbalance, collapses to unconditionally predicting the
majority class. Both the fine-tuning prompt template and completion format
exactly match variant V1's own zero-component prompt, so the fine-tuned
model is compared like-for-like against the prompting-only results at
inference time: same instruction, same input format, only the weights
differ. We save a per-epoch adapter snapshot and promote the epoch with
the highest F1 on the EDOS dev split to the one evaluated on the test
split, mirroring the encoder baseline's own selection criterion below.

\textbf{Encoder.} A \texttt{roberta-base} sequence classifier
(2 output labels), fine-tuned for 3 epochs (learning rate
$2{\times}10^{-5}$, batch size 64, max sequence length 128) directly on
the training split, with the epoch achieving the best dev-split F1
promoted to the reported test-set number.

\textbf{Classical.} A non-neural baseline: TF-IDF features
(unigrams and bigrams, minimum document frequency 2, 50,000-feature cap,
sublinear term-frequency scaling) feeding a logistic regression classifier
(balanced class weighting to offset EDOS's natural imbalance, up to 1,000
solver iterations), fit on the same training split.

\subsection{Statistical Testing}
\label{sec:methodology:stats}

For each model we test every non-baseline variant (V2--V18) against V1
with two independent, per-model-family-corrected tests: an exact
McNemar's test on the accuracy-based discordant-pair counts between the
two variants' predictions, and a paired bootstrap test on the difference
in F1 (2,000 resamples, percentile confidence interval). Both tests'
p-values are Holm-corrected within their own family of 17 comparisons per
model, rather than pooled across the two tests. McNemar and the F1-diff
bootstrap test the same baseline-vs-variant hypothesis via different
statistics, so correcting one family with the other's p-values would
misrepresent the family-wise error control each provides on its own
(Sec.~\ref{sec:results:significance} discusses a case where the two tests
disagree).

\subsection{Sensitivity-Sweep Design}
\label{sec:methodology:sweeps}

Three supplementary robustness checks target each model's own
best-performing main-grid variant only, not the full grid along a second
axis, which would be prohibitively expensive: (i) an \textbf{order sweep},
rerendering that variant's active components under up to 5 distinct,
seeded random orderings of \{context, chain-of-thought, aspects,
few-shot\} (deduplicated to distinct induced orders, since an inactive
component does not affect ordering), holding the active component set
itself fixed; (ii) a \textbf{few-shot-size sweep}, rerunning the variant
at $k\in\{1,3,4\}$ demonstrations per class against the grid's own default
$k{=}2$, for variants where few-shot is active; and (iii) a
\textbf{demonstration-seed sweep}, redrawing which $k{=}2$-per-class
demonstrations are sampled under three alternate random seeds (1, 7, 123)
against the grid's own fixed seed (42), holding both the component order
and the demonstration count fixed, for variants where few-shot is active.
All three sweeps reuse the same evaluation engine and generation settings
as the main grid (Sec.~\ref{sec:experimental:generation});
Sec.~\ref{sec:results:sensitivity} and
Sec.~\ref{sec:results:seed-sensitivity} report the results.


\section{Experimental Setup}
\label{sec:experimental}

\subsection{Dataset}
\label{sec:experimental:dataset}

We use the full official release of the Explainable Detection of Online
Sexism (EDOS) dataset from SemEval-2023 Task 10 \cite{kirk2023semeval}:
20,000 English-language social media posts labeled sexist / not sexist,
with every sexist post further annotated into one of four categories
(threats/plans to harm/incitement, derogation, animosity, and prejudiced
discussions). We use the dataset's own official train/dev/test partition
(14,000/2,000/4,000), fetched from a pinned upstream commit and verified
by SHA-256 checksum against the exact release used throughout this study,
so an upstream force-push cannot silently change what was evaluated.
Class balance is stable at roughly 24\% sexist / 76\% not sexist across
all three splits (train: 3,398 sexist / 10,602 not sexist; dev: 486 / 1,514;
test: 970 / 3,030, the exact figures underlying
Sec.~\ref{sec:results:overview}'s base rate).

The training split supplies both the few-shot demonstration pool
(Sec.~\ref{sec:methodology:prompts}) and the fine-tuning data for all
three baselines (Sec.~\ref{sec:methodology:baselines}); the dev split is
used only for the fine-tuned baselines' epoch selection
(Sec.~\ref{sec:methodology:baselines}). All reported metrics (prompting
grid, baselines, and sensitivity sweeps alike) are computed on the same
held-out test split. For the subtype-level error analysis
(Sec.~\ref{sec:results:subtype}), the test split's 970 sexist examples
break down as: derogation (454), animosity (333), prejudiced discussions
(94), and threats/plans to harm/incitement (89).

\subsection{Inference Precision}
\label{sec:experimental:precision}

Generation uses bf16 throughout, with no additional post-training
quantization: calibration showed 4-bit NF4 quantization and bf16 cost
effectively the same wall-clock time per example (30.28s vs.\ 30.47s), so
bf16 was kept as the default given ample VRAM headroom across the
roster, rather than paying the complexity of a quantized inference path
for no measured throughput benefit. This is separate from the QLoRA
baseline (Sec.~\ref{sec:methodology:baselines}), which always quantizes
its frozen base weights to 4-bit NF4 as standard QLoRA practice; there,
quantization is what makes fine-tuning a full-size 7--9B model on a
single GPU feasible at all, an entirely different concern from
generation-time speed.

\subsection{Generation and Decoding}
\label{sec:experimental:generation}

All generation is greedy (\texttt{do\_sample=False}, \texttt{top\_p=1.0},
no repetition penalty) under a fixed global seed (42) throughout.
Non-chain-of-thought variants use a 2-token generation budget, matching
their forced trailing \texttt{ANSWER:} cue
(Sec.~\ref{sec:methodology:prompts}); chain-of-thought variants use a
220-token budget, calibrated for the five non-reasoning models' brief
prose reasoning. qwen3 is the exception: its native thinking mode can
reason for far longer than 220 tokens before reaching a decision, so its
chain-of-thought units use a 1,024-token budget instead, directly
reflected in its substantially higher per-example latency
(Sec.~\ref{sec:results:cost}). Generation batch size is 32 for five of the
six models, calibrated to stay within available GPU memory (zero
out-of-memory failures observed at that size); qwen3 uses batch size
128, calibrated the same way for its longer generations
(fail\_ratio 0.008 at that batch size).

\subsection{Evaluation Metrics}
\label{sec:experimental:metrics}

We report accuracy and binary precision/recall/F1 with the sexist class as
positive, computed identically for every (model, variant) unit, baseline,
and sweep condition. A response is parsed by locating its \emph{last}
occurrence of an answer marker (\texttt{ANSWER:}, or a residual
chat-template artifact) and reading the YES/NO token that follows; this
single rule handles both the prompt-inserted marker of non-chain-of-thought
variants and the model-produced marker at the end of a chain-of-thought
response without special-casing either. A response that never reaches an
unambiguous YES/NO decision (including the deliberately unresolved case
of both words appearing with no leading marker) is a parse failure; the
fraction of failures is reported per unit as \texttt{fail\_ratio} and, in
the headline metrics used throughout this paper, failures default to a
``not sexist'' prediction. This default is a real cost only when
\texttt{fail\_ratio} is non-trivial, which is itself informative
(Sec.~\ref{sec:results:llama-refusal}).


\section{Results and Analysis}
\label{sec:results}

\subsection{Experimental Overview}
\label{sec:results:overview}

All results below are computed against the full official EDOS test split
\cite{kirk2023semeval}, $n=4{,}000$, at its real class distribution:
3{,}030 not-sexist examples (75.75\%) and 970 sexist examples (24.25\%).
We state this base rate up front because several of the accuracy figures
reported in this section are only interpretable against it: a
classifier that predicts ``not sexist'' unconditionally would already
score 0.758 accuracy by construction.

\subsection{Main Prompt-Design Grid}
\label{sec:results:grid}

Table~\ref{tab:main-grid} reports F1 for all 18 prompt variants across
the six evaluated models. Bold marks the best model for a given variant
(row); underline marks each model's own best-performing variant
(column).

\begin{table*}[t]
\centering
\footnotesize
\caption{F1 score by prompt variant (V1--V18) and model, full EDOS test
split ($n=4{,}000$). R = role, A = aspects, C = context, T =
chain-of-thought, F = few-shot ($k{=}2$ per class). Bold: best model per
row. Underline: best variant per model (column).
\textsuperscript{$\dagger$}llama3.1's V14 is omitted: 54.6\% of its
generations are safety refusals rather than classification attempts
(Sec.~\ref{sec:results:llama-refusal}). llama3.1's V8, V13, and V15
cells (F1 0.240, 0.361, 0.360) are reported but carry a related caveat:
39--49\% of their generations also fail to answer, mostly by the same
refusal mechanism (Table~\ref{tab:llama-cot}).}
\label{tab:main-grid}
\setlength{\tabcolsep}{4pt}
\begin{tabular}{l ccccc cccccc}
\toprule
& \multicolumn{5}{c}{Prompt Features} & \multicolumn{6}{c}{F1 Score} \\
\cmidrule(lr){2-6} \cmidrule(lr){7-12}
Variant & R & A & C & T & F & Mistral-7B & Phi-3-mini & Llama-3.1-8B & Qwen2.5-7B & Gemma-2-9B & Qwen3-8B \\
\midrule
V1 & -- & -- & -- & -- & -- & 0.467 & 0.138 & 0.441 & \textbf{0.537} & 0.448 & 0.481 \\
V2 & \checkmark & -- & -- & -- & -- & 0.454 & 0.256 & 0.472 & \textbf{0.535} & 0.444 & 0.483 \\
V3 & -- & \checkmark & -- & -- & -- & 0.490 & 0.413 & 0.550 & \textbf{0.570} & 0.501 & 0.543 \\
V4 & -- & -- & \checkmark & -- & -- & 0.453 & 0.056 & 0.466 & \textbf{0.493} & 0.454 & 0.477 \\
V5 & -- & -- & -- & \checkmark & -- & 0.496 & 0.480 & 0.412 & 0.509 & 0.443 & \textbf{0.553} \\
V6 & -- & -- & -- & -- & \checkmark & 0.567 & 0.517 & \underline{0.550} & \underline{0.585} & 0.530 & \textbf{0.590} \\
V7 & \checkmark & -- & -- & -- & \checkmark & \underline{0.567} & 0.533 & 0.519 & 0.583 & \underline{0.530} & \textbf{0.594} \\
V8 & -- & -- & -- & \checkmark & \checkmark & 0.471 & 0.438 & 0.240 & 0.571 & 0.503 & \underline{\textbf{0.627}} \\
V9 & \checkmark & \checkmark & -- & -- & \checkmark & 0.559 & \underline{\textbf{0.566}} & 0.471 & 0.559 & 0.514 & 0.562 \\
V10 & \checkmark & \checkmark & -- & \checkmark & -- & 0.461 & 0.478 & 0.497 & 0.543 & 0.454 & \textbf{0.544} \\
V11 & -- & \checkmark & \checkmark & \checkmark & -- & 0.483 & 0.466 & 0.508 & 0.556 & 0.471 & \textbf{0.559} \\
V12 & -- & \checkmark & -- & \checkmark & \checkmark & 0.462 & 0.461 & 0.501 & 0.568 & 0.496 & \textbf{0.602} \\
V13 & \checkmark & -- & -- & \checkmark & \checkmark & 0.473 & 0.448 & 0.361 & 0.576 & 0.482 & \textbf{0.614} \\
V14 & -- & -- & \checkmark & \checkmark & \checkmark & 0.447 & 0.461 & \textsuperscript{$\dagger$}-- & 0.571 & 0.500 & \textbf{0.618} \\
V15 & \checkmark & -- & \checkmark & \checkmark & \checkmark & 0.452 & 0.464 & 0.360 & 0.581 & 0.479 & \textbf{0.618} \\
V16 & \checkmark & \checkmark & \checkmark & \checkmark & -- & 0.474 & 0.473 & 0.495 & 0.537 & 0.454 & \textbf{0.546} \\
V17 & \checkmark & \checkmark & -- & \checkmark & \checkmark & 0.440 & 0.486 & 0.470 & 0.556 & 0.485 & \textbf{0.586} \\
V18 & \checkmark & \checkmark & \checkmark & \checkmark & \checkmark & 0.439 & 0.489 & 0.494 & 0.572 & 0.477 & \textbf{0.591} \\
\bottomrule
\end{tabular}
\end{table*}

Two patterns hold across all six models. First, the zero-shot baseline
(V1) is never any model's best variant, and the fullest five-component
prompt (V18) is not either: for every model, the top score comes from
a two- or three-component combination that includes few-shot examples.
Second, model behavior at V1 varies far more than the eventual ceiling
does: qwen2.5 opens at F1~0.537 with no prompt engineering at all, while
phi3 opens at F1~0.138, driven by a recall of just 0.082. It flags
almost nothing as sexist until given role framing, aspect decomposition,
or examples to work from. Qwen2.5 is the strongest model for every
V1--V4 (role/aspects/context-only) variant, but qwen3 becomes the
strongest model for every variant from V5 onward, once chain-of-thought
or few-shot enters the prompt. Qwen3's own best variant, V8
(chain-of-thought plus few-shot, F1~0.627), is also the single highest
score in the entire grid.

\subsection{Qualitative Finding: Safety-Alignment Interference in Llama-3.1-8B}
\label{sec:results:llama-refusal}

Llama-3.1-8B's V14 (context, chain-of-thought, few-shot; no role or
aspect framing) is omitted from Table~\ref{tab:main-grid} rather than
scored: 54.6\% of its 4{,}000 generations do not reach a parseable
YES/NO decision, reproduced identically across two independent runs
under this study's deterministic (greedy) decoding. Reading the raw
generations, these are not truncated or malformed outputs: they are
explicit safety refusals (e.g., ``I cannot create content that is
discriminatory or hateful''). Chain-of-thought is the trigger: no
variant without it shows a meaningful \texttt{fail\_ratio} for this
model. V14 is the worst single case of a pattern that runs across every
chain-of-thought variant llama3.1 was evaluated on, detailed below, not
an isolated one.

Scoring refusals as ``not sexist,'' the standard convention this study
uses elsewhere, V14 reaches 0.726 accuracy (almost exactly the test
split's 0.758 not-sexist base rate), with recall collapsing to 0.278
and F1 at 0.330. Restricted to the 45.4\% of rows the model actually
answers, F1 recovers to 0.525, in line with llama3.1's better-scoring
variants (V3 and V6, both F1~0.550). The two views tell different
stories: the model does not appear worse at the task when it engages,
but it declines to engage on over half the inputs under this specific
prompt structure. The answered subset is not a random sample either:
refusal plausibly correlates with how explicit the source text is. This
is a finding about llama3.1's safety alignment interacting with prompt
structure, not a gap in the evaluation; full derivation and the two-run
reproduction are documented alongside the raw refusal transcripts in the
project's supplementary material.

\begin{table}[t]
\centering
\footnotesize
\caption{\texttt{fail\_ratio} for llama3.1's ten other chain-of-thought
variants (none excluded from Table~\ref{tab:main-grid}; all stay below
the 0.5 exclusion gate). Refusal share is the fraction of a variant's
failed rows matching an explicit decline pattern (e.g.\ ``I cannot,''
``not able to''); the remainder are token-budget truncations, not
declines (see text).}
\label{tab:llama-cot}
\begin{tabular}{lcccc}
\toprule
Variant & Components & Aspects? & fail\_ratio & Refusal share \\
\midrule
V8  & T+F       & --         & 0.491 & 0.98 \\
V15 & R+C+T+F   & --         & 0.441 & 0.94 \\
V13 & R+T+F     & --         & 0.395 & 0.90 \\
V5  & T         & --         & 0.236 & 0.82 \\
V11 & A+C+T     & \checkmark & 0.214 & 0.31 \\
V12 & A+T+F     & \checkmark & 0.214 & 0.54 \\
V17 & R+A+T+F   & \checkmark & 0.171 & 0.63 \\
V16 & R+A+C+T   & \checkmark & 0.166 & 0.46 \\
V18 & R+A+C+T+F & \checkmark & 0.146 & 0.32 \\
V10 & R+A+T     & \checkmark & 0.121 & 0.45 \\
\bottomrule
\end{tabular}
\end{table}

Table~\ref{tab:llama-cot} reports \texttt{fail\_ratio} for llama3.1's
ten other chain-of-thought variants. Classifying each variant's failed
rows by whether they match an explicit refusal pattern reveals two
distinct causes, not one, and which dominates splits cleanly on whether
aspects is active. Every variant without aspects is refusal-dominant
(82--98\% of its failures are explicit declines, averaging around 100
characters each, matching V14's own transcripts). Every variant with
aspects is truncation-dominant or close to an even split (31--63\%
refusal); its non-refusal failures average 1{,}000--1{,}080 characters,
sitting right at this study's 220-token chain-of-thought budget
(Sec.~\ref{sec:experimental:generation}), and on inspection consist of
genuine step-by-step reasoning cut off mid-sentence before reaching a
decision, not a truncated refusal message. Aspects asks the model to
work through seven named linguistic markers before answering; for
llama3.1, that reasoning routinely does not finish inside a token budget
calibrated for the other five models' shorter prose reasoning. V8
(\texttt{fail\_ratio} 0.491), V13 (0.395), and V15 (0.441) are the three
most refusal-affected cells still reported in Table~\ref{tab:main-grid}
(F1 0.240, 0.361, and 0.360 respectively); their scores carry the same
caveat as V14's, just below rather than above the exclusion threshold: a
property of this specific model, worth knowing before comparing its
numbers to the other five.

\subsection{Statistical Significance of Prompt-Design Effects}
\label{sec:results:significance}

For each model we test every non-baseline variant against V1 with
exact McNemar's test (on the discordant-pair counts between the two
variants' predictions) and a paired bootstrap test on the F1 difference
(2{,}000 resamples), both Holm-corrected across the model's variants.
Table~\ref{tab:significance} summarizes the results.

\begin{table}[t]
\centering
\footnotesize
\caption{Significance of prompt-design effects vs.\ baseline V1,
Holm-corrected within each model ($\alpha=0.05$). ``Compared'' excludes
V1 itself and, for llama3.1, the omitted V14. ``McN.\ sig.''/``F1 sig.''
are the counts of variants significant under exact McNemar's test and
the paired F1-diff bootstrap, respectively (Sec.~\ref{sec:methodology:stats}).}
\label{tab:significance}
\setlength{\tabcolsep}{4pt}
\begin{tabular}{lccccc}
\toprule
Model & Compared & McN. sig. & F1 sig. & Best var. & Best $\Delta$F1 \\
\midrule
Mistral-7B & 17 & 17 & 12 & V7 & +0.100 \\
Phi-3-mini & 17 & 13 & 17 & V9 & +0.428 \\
Llama-3.1-8B & 16 & 16 & 16 & V6 & +0.110 \\
Qwen2.5-7B & 17 & 14 & 14 & V6 & +0.049 \\
Gemma-2-9B & 17 & 16 & 14 & V7 & +0.082 \\
Qwen3-8B & 17 & 16 & 15 & V8 & +0.146 \\
\midrule
Total & 101 & 92 & 88 & -- & -- \\
\bottomrule
\end{tabular}
\end{table}

88 of the 101 variant-vs-baseline comparisons across all six models
show a significant F1 difference, and 92 of 101 show a significant
McNemar effect: the large majority of this study's prompt-design choices
produce a real, statistically distinguishable effect on model behavior,
not sampling noise. The two tests disagree on a substantial fraction of
comparisons, in both directions (mistral: 17 McNemar-significant vs.\ 12
F1-diff-significant; phi3: 13 vs.\ 17). This divergence is expected,
since the two tests measure different quantities; phi3's V1$\to$V9 jump,
below, shows why.

Phi-3-mini's V1$\to$V9 jump is the clearest illustration. F1 rises from
0.138 to 0.566 (+0.428, by far the largest single effect in this study),
while accuracy moves only from 0.750 to 0.759 and the McNemar test does
not reach significance after correction ($p_{\text{holm}}=0.695$). The
reason is arithmetic, not a contradiction between tests: at V1, recall
is 0.082 and precision 0.421, implying on the order of 80 true positives
against 110 false positives out of 4{,}000 predictions. At V9, recall
reaches 0.647 and precision 0.502, implying roughly 628 true positives
against 622 false positives. The variant gains around 548 correct
sexist-flags and loses around 512 new false alarms, almost, but not
exactly, offsetting on a 4{,}000-row test set that is 76\% negative, so
raw accuracy (and the McNemar test built on it) barely moves. F1, which
is computed against the 970-row positive class rather than the full
test set, is not diluted by that majority-class denominator, so it
reflects the same underlying shift as a large, clearly significant
change. This is the concrete case for reporting both test families
rather than collapsing them into one number, and it is the opposite
pattern from every other model in this study, where McNemar catches
more significant comparisons than the F1-diff bootstrap does.

\subsection{Comparison Against Fine-Tuned and Classical Baselines}
\label{sec:results:baselines}

Table~\ref{tab:baselines} places each model's best training-free prompt
variant against a QLoRA-tuned version of the same model, a single
TF-IDF plus logistic-regression classifier (no LLM), and a
fine-tuned RoBERTa-base encoder.

\begin{table}[t]
\centering
\footnotesize
\caption{Best training-free prompting vs.\ fine-tuned and classical
baselines. Classical (TF-IDF + logistic regression): F1 0.617, accuracy
0.799. Encoder (RoBERTa-base): F1 0.735, accuracy 0.874. $\Delta$
columns are F1 relative to the classical baseline.}
\label{tab:baselines}
\begin{tabular}{lccc}
\toprule
Model & Best prompt F1 & QLoRA F1 & $\Delta$ vs.\ classical \\
\midrule
Mistral-7B & 0.567 (V7) & 0.788 & $-$0.050 \\
Phi-3-mini & 0.566 (V9) & 0.734 & $-$0.051 \\
Llama-3.1-8B & 0.550 (V6) & 0.795 & $-$0.067 \\
Qwen2.5-7B & 0.585 (V6) & 0.786 & $-$0.032 \\
Gemma-2-9B & 0.530 (V7) & 0.794 & $-$0.087 \\
Qwen3-8B & 0.627 (V8) & 0.779 & $+$0.010 \\
\bottomrule
\end{tabular}
\end{table}

The headline result of this comparison is not favorable to prompting.
\textbf{Five of the six models' best training-free prompt variant, after
an 18-point search over prompt design, scores below a bag-of-words
logistic regression classifier}: a baseline with no language model,
no GPU, and a fraction of a second of CPU inference per example. Only
qwen3 edges past it, by F1~+0.010, and even qwen3 loses to it on
accuracy (0.773 vs.\ 0.799). QLoRA fine-tuning and the RoBERTa encoder
are in a different tier entirely: every QLoRA-tuned model reaches
F1~0.734--0.795, and the encoder reaches F1~0.735, both comfortably
above the classical baseline and far above every prompted number.
Ranked by F1, the ordering is consistent across the whole study:
QLoRA~$\approx$~encoder~$\gg$~classical~$>$~training-free prompting,
with qwen3 as the one partial exception at the classical boundary. This
is the direct answer to this study's ``no non-LLM baseline'' gap: the
classical comparison reframes what the prompt-design results in
Sec.~\ref{sec:results:grid}--\ref{sec:results:significance} mean.
Prompt engineering produces large, statistically real differences
\emph{within} the space of training-free approaches, but that entire
space sits, for five of six models, below what a linear classifier over
word counts already achieves on this task.

\subsection{Error Analysis by Sexism Subtype}
\label{sec:results:subtype}

Table~\ref{tab:subtype} breaks down each model's best variant by the
four fine-grained sexism categories in the EDOS taxonomy
\cite{kirk2023semeval}, reporting the false-negative rate (share of
truly sexist examples in that category the model misses) for each.

\begin{table}[t]
\centering
\footnotesize
\caption{False-negative rate by sexism subtype, best variant per model.
1: threats/plans to harm/incitement ($n=89$). 2: derogation ($n=454$).
3: animosity ($n=333$). 4: prejudiced discussions ($n=94$).}
\label{tab:subtype}
\begin{tabular}{lcccccc}
\toprule
Model & Variant & 1 & 2 & 3 & 4 & Mean \\
\midrule
Mistral-7B & V7 & 0.011 & 0.057 & 0.147 & 0.181 & 0.099 \\
Phi-3-mini & V9 & 0.337 & 0.300 & 0.372 & 0.553 & 0.391 \\
Llama-3.1-8B & V6 & 0.079 & 0.064 & 0.129 & 0.255 & 0.132 \\
Qwen2.5-7B & V6 & 0.281 & 0.141 & 0.303 & 0.309 & 0.258 \\
Gemma-2-9B & V7 & 0.022 & 0.015 & 0.048 & 0.043 & 0.032 \\
Qwen3-8B & V8 & 0.303 & 0.150 & 0.237 & 0.372 & 0.266 \\
\bottomrule
\end{tabular}
\end{table}

Even at each model's own best variant, subtype-level performance
diverges far more than the aggregate F1 numbers in
Table~\ref{tab:main-grid} suggest. Gemma-2-9B misses under 5\% of
threats and derogation cases; phi3, at its own best variant, still
misses 55\% of ``prejudiced discussions'' and over a third of every
other category. Category 4 (prejudiced discussions, the smallest class
at $n=94$) is the hardest subtype for five of the six models. This
gap widens sharply when the worst rather than the best variant is
considered: phi3's worst variant (V4, context-only) misses 96--100\% of
sexist examples in every subtype, a near-total collapse into
predicting ``not sexist,'' while gemma2's worst variant (V8) never
exceeds a 16\% false-negative rate in any subtype. The spread between a
model's best-case and worst-case prompt is itself a per-model property,
not a constant: gemma2 is the most robust to a bad prompt choice in this
study, and phi3 the least.

\subsection{Sensitivity to Prompt-Component Order and Few-Shot Size}
\label{sec:results:sensitivity}

For each model's best main-grid variant, we re-evaluate under (a) an
alternate ordering of its prompt components, where the variant has more
than one component the harness can meaningfully reorder, and (b) an
alternate few-shot count ($k \in \{1,3,4\}$ per class, against the
main grid's default $k=2$). Table~\ref{tab:sensitivity} reports each as
an F1 delta from the main-grid value.

\begin{table}[t]
\centering
\footnotesize
\caption{Sensitivity of each model's best variant to component reordering
and few-shot count, as $\Delta$F1 from the Table~\ref{tab:main-grid}
value. n/a: variant not eligible for order sweep (see text).}
\label{tab:sensitivity}
\begin{tabular}{lcccc}
\toprule
Model (variant) & Order & $k{=}1$ & $k{=}3$ & $k{=}4$ \\
\midrule
Mistral-7B (V7) & n/a & $-$0.039 & +0.024 & +0.028 \\
Phi-3-mini (V9) & $-$0.032 & +0.024 & $-$0.019 & $-$0.092 \\
Llama-3.1-8B (V6) & n/a & +0.023 & $-$0.004 & +0.001 \\
Qwen2.5-7B (V6) & n/a & $-$0.031 & $-$0.005 & $-$0.002 \\
Gemma-2-9B (V7) & n/a & $-$0.018 & $-$0.011 & +0.011 \\
Qwen3-8B (V8) & +0.004 & $-$0.005 & $-$0.008 & +0.018 \\
\bottomrule
\end{tabular}
\end{table}

Order sensitivity applies only to phi3 (V9: role, aspects, few-shot) and
qwen3 (V8: chain-of-thought, few-shot). The other four models' best
variants either combine a single component with few-shot or combine two
components that render identically regardless of order under this
study's prompt template, so no meaningful reordering exists to test.
Where it does apply, the two models diverge sharply: qwen3 is stable
under reordering ($\Delta=+0.004$), while phi3 loses 0.032 F1 under the
alternate order, a real, if modest, position effect for this model
specifically.

The few-shot-size sweep carries a more general caveat. This study's
main-grid results all use $k=2$ examples per class, but that is not a
tuned optimum: four of six models score higher F1 at $k=4$ than at the
$k=2$ default (mistral +0.028, gemma2 +0.011, qwen3 +0.018, llama3.1
+0.001), meaning the headline numbers in Table~\ref{tab:main-grid}
understate what few-shot prompting can achieve for most of this study's
models. Phi-3-mini is the exception, and a stark one: F1 falls
monotonically as $k$ increases past the default, down to $-0.092$ at
$k=4$, driven by precision rising (0.530$\to$0.670) as recall collapses
(0.663$\to$0.366). Where the other five models either improve or stay
roughly flat with more examples, phi3 becomes systematically more
conservative: the same under-triggering behavior visible in its V1
baseline (Sec.~\ref{sec:results:grid}) re-emerges as few-shot count
grows, rather than being fixed by it.

\subsection{Sensitivity to Demonstration Selection}
\label{sec:results:seed-sensitivity}

The order and size sweeps above vary how the few-shot demonstrations are
arranged and how many are shown, but never which specific examples are
drawn: every main-grid variant samples its $k{=}2$-per-class
demonstrations once, with a fixed seed (42). Table~\ref{tab:seed-sensitivity}
reruns each model's best variant under three alternate seeds (1, 7, 123),
reporting the resulting F1 range as a delta from the main-grid value.

\begin{table}[t]
\centering
\footnotesize
\caption{Sensitivity of each model's best variant to which few-shot
demonstrations are drawn: F1 range across 3 alternate seeds (1, 7, 123),
as $\Delta$F1 from the Table~\ref{tab:main-grid} value (grid's own seed,
42).}
\label{tab:seed-sensitivity}
\begin{tabular}{lcc}
\toprule
Model (variant) & Grid F1 & $\Delta$F1 range \\
\midrule
Mistral-7B (V7) & 0.567 & $-$0.025 to +0.011 \\
Phi-3-mini (V9) & 0.566 & $-$0.235 to +0.007 \\
Llama-3.1-8B (V6) & 0.550 & $-$0.097 to $-$0.028 \\
Qwen2.5-7B (V6) & 0.585 & $-$0.037 to $-$0.011 \\
Gemma-2-9B (V7) & 0.530 & $-$0.034 to $-$0.012 \\
Qwen3-8B (V8) & 0.627 & $-$0.025 to $-$0.003 \\
\bottomrule
\end{tabular}
\end{table}

Four of six models stay within a narrow, consistently negative range
(gemma2, qwen2.5, and qwen3 within 0.02--0.04 F1; llama3.1 within
0.03--0.10 F1): every alternate seed tested scores at or below the main
grid's own draw. Phi-3-mini is again this study's outlier, and by the
widest margin observed in any sensitivity check reported here: F1 falls
from 0.566 to 0.330 under one alternate seed. This is not a parsing
failure, \texttt{fail\_ratio} is 0 on every seed tested, main grid
included, but a genuine behavioral shift: precision rises to 0.744 while
recall collapses to 0.212, the same conservative, under-triggering
pattern already observed as $k$ grows past the default
(Sec.~\ref{sec:results:sensitivity}) and in this model's V1 baseline
(Sec.~\ref{sec:results:grid}), now shown to also depend on which four
examples happen to be drawn. The main grid's own seed scores at or above
all three alternate seeds for four of six models (llama3.1, qwen2.5,
gemma2, qwen3); the two exceptions are mistral, where seeds 1 and 7 each
score marginally higher (+0.011, +0.002 F1), and phi3, where seed 7
scores 0.007 F1 higher despite this study's largest sensitivity-sweep
swing in the opposite direction coming from a different seed on the same
model. Seed 42 was fixed as an arbitrary default before this sweep was
run, not chosen for its result, but the pattern means
Table~\ref{tab:main-grid}'s headline numbers sit, for most models, closer
to the top than the middle of what a different demonstration draw would
produce.

\subsection{Computational Cost}
\label{sec:results:cost}

Table~\ref{tab:cost} reports wall-clock inference cost for the full
18-variant grid per model, recorded automatically as part of every run.

\begin{table}[t]
\centering
\footnotesize
\caption{Inference cost across the full 18-variant grid, per model.}
\label{tab:cost}
\begin{tabular}{lccc}
\toprule
Model & Total (h) & Mean (s/ex.) & Slowest variant \\
\midrule
Mistral-7B & 3.92 & 0.196 & V10, 25.9 min \\
Phi-3-mini & 3.19 & 0.160 & V12, 24.0 min \\
Llama-3.1-8B & 4.71 & 0.249 & V18, 31.5 min \\
Qwen2.5-7B & 1.51 & 0.075 & V11, 14.0 min \\
Gemma-2-9B & 6.26 & 0.313 & V18, 38.5 min \\
Qwen3-8B & 23.92 & 1.196 & V18, 146.8 min \\
\bottomrule
\end{tabular}
\end{table}

Qwen3-8B is an outlier by a wide margin: its full grid takes roughly
5--16$\times$ longer than any other model's, and its mean per-example
cost (1.196s) is 4--16$\times$ every other model's. This is a direct
consequence of its native reasoning mode, which this study leaves
enabled for chain-of-thought variants rather than capping its token
budget to match the other five models. Qwen3 is also this study's
best-performing model overall (Table~\ref{tab:main-grid},
\ref{tab:significance}); that performance is not free, and the
accuracy/cost trade-off it represents is a separate question from
whether it is the strongest model on accuracy alone.

\subsection{Confidence-Based AUC: A Second Comparison Against the Classical Baseline}
\label{sec:results:auc}

For the seven non-chain-of-thought variants (V1--V4, V6, V7, V9), the
model's decision is its first generated token, so a calibrated
P(sexist) score can be read directly off that token's YES/NO logits,
softmax-normalized between the two candidates, without generating
anything further. Chain-of-thought variants are excluded: their
decision only emerges after free-form reasoning, so no single generated
token stands in for it. We extract this score for all six models across
the full 4,000-example test split and score it against the true labels
with ROC-AUC. As a validity check, the confidence-implied call
(P(sexist) $\geq$ 0.5) agrees with the model's actual generated decision
on 94.5--99.7\% of rows for every model's variant in
Table~\ref{tab:auc}, confirming the first-token assumption holds on
real output, not only by construction.

\begin{table}[t]
\centering
\footnotesize
\caption{AUC from first-token P(sexist), best non-chain-of-thought
variant per model. Classical baseline (TF-IDF + logistic regression,
\texttt{predict\_proba}): AUC 0.848. $\Delta$ is AUC relative to the
classical baseline. \textsuperscript{$\dagger$}Qwen3-8B's true best
variant, V8 (F1 0.627, Table~\ref{tab:baselines}), is chain-of-thought
and has no single decision token this method can score; V7 is its best
scorable stand-in, not its true best variant.}
\label{tab:auc}
\begin{tabular}{lcccc}
\toprule
Model & Variant & F1 & AUC & $\Delta$ vs.\ classical \\
\midrule
Mistral-7B & V7 & 0.567 & 0.843 & $-$0.005 \\
Phi-3-mini & V9 & 0.566 & 0.810 & $-$0.038 \\
Llama-3.1-8B & V6 & 0.550 & 0.832 & $-$0.016 \\
Qwen2.5-7B & V6 & 0.585 & 0.812 & $-$0.036 \\
Gemma-2-9B & V7 & 0.530 & 0.875 & $+$0.027 \\
Qwen3-8B\textsuperscript{$\dagger$} & V7 & 0.594 & 0.838 & $-$0.010 \\
\bottomrule
\end{tabular}
\end{table}

Table~\ref{tab:auc} complicates the story in Table~\ref{tab:baselines}
rather than repeating it. On F1, qwen3 is the only model to beat the
classical baseline, and it does so with V8, a chain-of-thought variant
this AUC method cannot score at all. On AUC, qwen3's best scorable
variant loses to the classical baseline (0.838 vs.\ 0.848), while
gemma2, whose F1 (0.530) trails the classical baseline by the widest
margin in this study, is the only model whose AUC (0.875) beats it.
Which model, if any, ``beats'' the classical baseline depends on which
metric answers the question. Training-free prompting still fails to
consistently clear the classical baseline under either metric; the
exception to that just changes identity depending on which metric is
used.

\subsection{Confidence-Based Abstention: Trading Coverage for Precision}
\label{sec:results:abstention}

The same first-token P(sexist) score behind Table~\ref{tab:auc} also
supports a simple abstention policy: instead of forcing a decision on
every example, withhold a prediction whenever the score falls within a
symmetric band around the 0.5 decision boundary, deferring the closest
calls rather than answering them. For each model's Table~\ref{tab:auc}
variant, Table~\ref{tab:abstention} reports the resulting coverage (the
fraction of the test split still answered) and F1 on that covered subset
at a $\pm0.2$ band, against the same variant's full-coverage F1
(band$=$0) for reference.

\begin{table}[t]
\centering
\footnotesize
\caption{Confidence-based abstention at a $\pm0.2$ band around
P(sexist)$=$0.5, same variant per model as Table~\ref{tab:auc}. Coverage
is the fraction of the 4{,}000-example test split still answered after
abstention; F1 is computed on that covered subset only; $\Delta$F1 is
relative to the same variant's full-coverage (no-abstention) F1.}
\label{tab:abstention}
\begin{tabular}{lcccc}
\toprule
Model & Variant & Coverage & F1 & $\Delta$F1 \\
\midrule
Mistral-7B & V7 & 0.941 & 0.585 & $+$0.015 \\
Phi-3-mini & V9 & 0.628 & 0.691 & $+$0.109 \\
Llama-3.1-8B & V6 & 0.809 & 0.601 & $+$0.049 \\
Qwen2.5-7B & V6 & 0.971 & 0.593 & $+$0.007 \\
Gemma-2-9B & V7 & 0.959 & 0.541 & $+$0.011 \\
Qwen3-8B & V7 & 0.962 & 0.602 & $+$0.008 \\
\bottomrule
\end{tabular}
\end{table}

F1 improves for every model at this band. Beyond being calibrated in
aggregate (Table~\ref{tab:auc}), the confidence score tracks
per-example difficulty: the examples it flags as uncertain really are,
on average, more error-prone. The size of that improvement, and the
coverage it costs to get it, is not uniform. Phi-3-mini gains the most
($+$0.109 F1) but only by abstaining on 37\% of the test split (coverage
0.628), a policy this study's other five models never need to reach a
comparable gain: mistral, gemma2, qwen2.5, and qwen3 all retain over
94\% coverage at the same band while gaining under 0.02 F1. Llama-3.1-8B
falls between the two patterns ($+$0.049 F1 at 81\% coverage). This
mirrors a pattern already visible elsewhere in this study: phi3's errors
are not spread evenly across the test set but concentrated in a
recoverable subset (Sec.~\ref{sec:results:sensitivity},
Sec.~\ref{sec:results:seed-sensitivity}, and
Sec.~\ref{sec:results:subtype} all single out phi3 as an outlier among
the six models), and abstention is one more lens that finds the same
thing. As with Table~\ref{tab:auc}, this policy is evaluated only on the
seven non-chain-of-thought variants; whether the same abstention gain
holds for chain-of-thought variants, where no single token yields a
confidence score, is untested by this method.


\section{Discussion}
\label{sec:discussion}

\subsection{What ``Prompt Design Matters'' Actually Means Here}
\label{sec:discussion:prompt-effects}

The statistical results (Sec.~\ref{sec:results:significance}) show a
real, measurable shift in model behavior for the large majority of this
study's prompt-design choices: phi3's role-plus-aspects-plus-few-shot
jump (+0.428 F1, V1$\to$V9) is the largest single effect this study's
significance tests measure. Fine-tuning was never put through the same
tests, but for this same model the QLoRA gain is larger still: F1 0.734
against V1's 0.138 (Table~\ref{tab:baselines}), a +0.596 improvement.
Prompt design's effect is large by the standards of training-free
methods; whether it is large by the standards of adapting the model's
own weights is a separate, smaller claim, and the baseline comparison
(Sec.~\ref{sec:results:baselines}) answers it: those effects operate
entirely within a space that, for five of six models, tops out below a
linear classifier trained on the same 14{,}000 examples in a fraction of
a second on a CPU. Prompt design is a real lever on this task, but it is
attached to a ceiling that a much simpler method already clears.

That ceiling moves, too. The few-shot-size sweep
(Sec.~\ref{sec:results:sensitivity}) shows four of six models score
higher at $k=4$ than the main grid's own $k=2$ default, so
Table~\ref{tab:main-grid} understates each model's real training-free
ceiling. Recomputing the baseline gap at each model's best available $k$
does not close it: mistral moves from $-0.050$ to $-0.022$ relative to
classical, gemma2 from $-0.087$ to $-0.076$, llama3.1 from $-0.067$ to
$-0.066$, all still below the classical baseline, while qwen3, already
ahead at $k=2$, extends its lead to $+0.028$. The direction of the
honest-limits finding survives its own most favorable adjustment; only
its margin moves.

One axis mattered less than prior in-context-learning work on other
tasks would predict. Lu et al.~\cite{lu2022fantastically} found that
reordering few-shot demonstrations alone can swing a model between
near-state-of-the-art and near-random performance. Reordering entire
prompt components produced nothing like that here: four of six models had
no meaningfully reorderable best variant at all, and where reordering did
apply, only phi3 showed a real effect ($-0.032$ F1). Component-order
stability is a genuine, if narrow, piece of good news for anyone
assembling a prompt in this space: which components to include is what
matters, not primarily what order to put them in.

\subsection{The Cost of the One Exception}
\label{sec:discussion:qwen3-cost}

qwen3 is the sole model whose best training-free prompt beats the
classical baseline on F1, and the paper's own framing needs to be honest
about what that win costs. Table~\ref{tab:cost} shows qwen3's full grid
running 5--16$\times$ longer than any other model's, driven by its native
reasoning mode left unconstrained on chain-of-thought variants
(Sec.~\ref{sec:experimental:generation}). Its margin over the classical
baseline is $+0.010$ F1; the classical baseline itself runs in a fraction
of a second per example with no GPU at all. A $+0.010$ F1 gain purchased
at 4--16$\times$ the inference cost of every other prompted model, and
still behind the classical baseline on accuracy (0.773 vs.\ 0.799), does
not read as an argument for training-free prompting over a linear
classifier in any setting where inference cost matters. Even the one
model that wins does not win by much, or cheaply.

That framing holds specifically under F1, the metric this study's main
grid and significance testing are built around; it does not hold under
every lens. Sec.~\ref{sec:results:auc}'s confidence-based AUC comparison
shows qwen3's best scorable variant losing to the classical baseline
(0.838 vs.\ 0.848), while gemma2, whose F1 (0.530) trails classical by
the widest margin of any model in this study, is the one model whose AUC
beats it. ``The one exception'' depends on which metric answers the
question, and qwen3's true best variant (V8) cannot be evaluated by the
AUC method at all. This does not weaken the cost argument above; it
means that argument is one honest answer among two.

\subsection{Safety Alignment as an Operational Risk}
\label{sec:discussion:refusal}

Llama-3.1-8B's V14 behavior (Sec.~\ref{sec:results:llama-refusal}) is
easy to misread as a parsing bug or a prompt-engineering mistake specific
to one variant. It is neither. The model reads a sexism-classification
request, once asked to reason step by step before answering, as a
request to generate hateful content rather than analyze it, and declines
on safety grounds over half the time in V14's specific configuration;
the same decline shows up, at lower but still substantial rates, across
every other chain-of-thought variant this model was evaluated on
(Table~\ref{tab:llama-cot}), so chain-of-thought reasoning itself, not
any one combination of surrounding components, is what triggers it.
Scored under this study's own default convention, that decline registers
as an unremarkable-looking F1 of 0.330 for V14, and quietly deflated
numbers for several other cells that stay in the main table, each
indistinguishable on its face from an ordinary weak variant. Only
inspecting the raw generations, rather than trusting the aggregate
metric, revealed what was actually happening; a related but mechanically
distinct pattern (Sec.~\ref{sec:results:llama-refusal}) shows that not
every elevated \texttt{fail\_ratio} in this model is a refusal at all,
some are a fixed reasoning-token budget proving too short once the
prompt also asks for aspect-by-aspect analysis, which is a generation
setting worth revisiting rather than a safety-alignment story. Röttger et
al.'s XSTest~\cite{rottger2023xstest} showed the exaggerated-safety
version of this failure mode under adversarial probing; llama3.1's
refusals show it can also emerge unprompted, from an ordinary
classification pipeline, on genuinely on-topic content. For any
deployment that scores a refusal as a default label rather than flagging
it, this is a silent failure mode: the system does not error, it
produces a plausible-looking, systematically wrong answer on a
substantial share of a specific input class.
\texttt{fail\_ratio} (Sec.~\ref{sec:experimental:metrics}), not accuracy
or F1 alone, is what surfaces it, and only because this study checked.

\subsection{What a Single Aggregate Score Hides}
\label{sec:discussion:aggregate}

The subtype breakdown and the significance-testing disagreement point at
the same underlying problem. The subtype breakdown
(Sec.~\ref{sec:results:subtype}) shows a model's own best variant can
still miss over half of the rarest sexism category, prejudiced
discussions, while its aggregate F1 looks unremarkable: gemma2 and
phi3 post worst-variant false-negative spreads of 16\% and 96--100\%
respectively, a gap the main grid's F1 column alone never reveals. The
significance-testing disagreement (Sec.~\ref{sec:results:significance})
reflects the same class-imbalance sensitivity: phi3's largest F1 effect
in the entire study is invisible to accuracy-based McNemar testing,
because a 76\%-negative test set can absorb a large shift in
positive-class detection without moving overall accuracy much at all.
Röttger et al.'s HateCheck~\cite{rottger2020hatecheck} made the general
version of this argument for hate speech detection: held-out accuracy
and F1 can look strong while hiding specific, systematic weaknesses that
only a finer-grained test reveals. This study's subtype table and its
dual significance tests apply that same argument to sexism detection, in
two independent places rather than one.

Together, these four points sharpen this study's central claim: prompt
design produces a genuine, well-measured effect on this task, but one
that operates inside a ceiling a simple non-LLM baseline already
exceeds for most models tested; the one model that clears that ceiling
pays a steep,
possibly disqualifying compute cost to do so; and the aggregate numbers
behind all of these conclusions can themselves hide model-specific
failure modes, a safety refusal or a rare-subtype blind spot, that only
closer inspection surfaces.


\section{Limitations}
\label{sec:limitations}

This study covers a single dataset, language, and pair of source
platforms: EDOS's English-language Reddit and Gab posts
\cite{kirk2023semeval}. Sexist language is expressed differently across
languages and platforms (Sec.~\ref{sec:related:datasets} discusses SWSR
\cite{jiang2021swsr}, the closest non-English counterpart), so nothing
here establishes that the same prompt-design effects, the same baseline
gap, or the same safety-refusal pattern would reproduce on a different
dataset. The study also targets only EDOS's binary sexist/not-sexist
subtask directly; the dataset's finer four-category taxonomy is used
solely for the post-hoc subtype error analysis
(Sec.~\ref{sec:results:subtype}), never as a training or prediction
target in its own right.

The prompt-design grid is a curated 18 of 32 possible component
combinations (Sec.~\ref{sec:methodology:prompts}), not an exhaustive
search, and every component's wording is hand-designed rather than
automatically optimized or drawn from a cited theoretical taxonomy; this
includes the \textbf{aspects} component's seven linguistic markers, which
reflect our own judgment of what sexist language looks like rather
than an established sociolinguistic framework. A negative result
within this space, most
prompted variants scoring below the classical baseline, is evidence
about the designs actually tested, not proof that no possible prompt
formulation could close the gap. Similarly, the few-shot demonstrations
for every main-grid variant are drawn once, with a fixed seed, from the
training split. Sec.~\ref{sec:methodology:sweeps}'s order, count, and
demonstration-seed sweeps test how sensitive each model's own best
variant is to arrangement, count, and which specific examples are drawn,
respectively (Sec.~\ref{sec:results:sensitivity},
Sec.~\ref{sec:results:seed-sensitivity}), but only for that one variant
per model, not the full grid; the seed sweep itself covers three
alternate draws per model, enough to surface phi3's large,
model-specific sensitivity to selection, but not an exhaustive
characterization of the selection space.

The model roster is six open-weight, instruction-tuned models in a
3.8--9B parameter range (Table~\ref{tab:models}). No larger open model
and no closed, commercial API model, the kind increasingly used in
production moderation systems, is evaluated, so neither the
prompt-design findings nor the baseline comparison necessarily extend to
that different and increasingly common deployment setting. Every
inference-cost figure in this study (Table~\ref{tab:cost}) is specific
to this study's own hardware and batching choices;
the direction of the cost argument in
Sec.~\ref{sec:discussion:qwen3-cost}, that qwen3's win over the
classical baseline is expensive, is likely to hold generally, but its
exact multiplier is not.

This study compares against a classical machine-learning baseline and
two fine-tuned neural baselines, but not against human inter-annotator
agreement on the same test items. Without that reference point, it is
not possible to say how close, or how far, any of this study's numbers
sit from whatever ceiling human disagreement on sexism labeling itself
imposes.

Sec.~\ref{sec:results:llama-refusal} and Table~\ref{tab:llama-cot} show
that llama3.1's elevated \texttt{fail\_ratio} on chain-of-thought
variants splits into two distinct causes depending on whether aspects is
active: explicit safety refusal when it is not, and token-budget
truncation, the model's own step-by-step reasoning running past this
study's fixed 220-token chain-of-thought budget, when it is. That budget
was calibrated for ``brief prose reasoning'' across the five
non-reasoning models generally
(Sec.~\ref{sec:experimental:generation}); it is demonstrably too short
for at least this one model under at least this one prompt
configuration. Whether the same refusal mechanism, at any scale,
generalizes to the other five models has since been checked
systematically rather than spot-checked: every chain-of-thought cell
with at least one failure across mistral, phi3, qwen2.5, gemma2, and
qwen3 (47 cells) was classified by the same refusal-pattern rule used
for Table~\ref{tab:llama-cot}, reading explicit-refusal
\emph{prevalence} as a share of the full 4{,}000-example test split
rather than a share of each cell's own failures, since the two can
diverge sharply when \texttt{fail\_ratio} itself is small. On that
measure, explicit refusal never exceeds 1.05\% of the test split for any
of the five (phi3's V13, the highest rate observed outside llama3.1;
mistral peaks at 0.30\%, gemma2 at 0.23\%, qwen3 at 0.15\%, and qwen2.5
shows none). Every one of llama3.1's own ten other chain-of-thought
cells sits well above that entire range, from 4.82\% at its lowest to
48.1\% at V8. Llama3.1's safety-refusal interference
(Sec.~\ref{sec:results:llama-refusal}) is, on this evidence, a property
of that model specifically, not a shared chain-of-thought failure mode
this study's prompt design triggers generally.

Sec.~\ref{sec:results:auc}'s AUC comparison, and
Sec.~\ref{sec:results:abstention}'s confidence-based abstention policy
built on the same score, both cover only the seven non-chain-of-thought
variants. The eleven chain-of-thought variants, including qwen3's own
best-F1 variant (V8), have no single decision token this method can
score, so neither AUC nor an abstention gain exists for them at all. The
abstention gain itself is also not free or uniform: Sec.~\ref{sec:results:abstention}
shows one model (phi3) buying a large F1 improvement only by abstaining
on over a third of the test split, while four of the other five models
gain comparatively little at a much smaller coverage cost, so the same
band width does not represent the same real-world tradeoff across
models.

Generation throughout this study is greedy and seeded, which removes
sampling randomness but not every source of nondeterminism: kernel-level
differences across GPU architectures, driver versions, and library
versions can still produce different outputs from bit-identical inputs
on different hardware. Each of five models' best variant was rerun on two cluster
nodes with the same GPU architecture (GA100, compute capability 8.0) but
different memory configurations, a 40GB-per-card node and an
80GB-per-card node: every one of 4{,}000 test examples produced
byte-identical raw output and an identical parsed decision on both
nodes, for all five models, confirming this study's pipeline is
bit-for-bit deterministic across nodes and memory configurations within
one GPU architecture. The two-run reproduction of llama3.1's V14 refusal
rate (Sec.~\ref{sec:results:llama-refusal}) is a special case of the
same result. qwen3, whose batch-128 footprint does not fit the 40GB
node's cards, was instead rerun on the 80GB node's A100 cards, again
matching Table~\ref{tab:main-grid}'s own V8 figure exactly. What this
study can report is strong, directly tested evidence for bit-identical
determinism across nodes and memory configurations within one GPU
architecture; whether the same holds across different GPU architectures
remains untested.


\section{Conclusion}
\label{sec:conclusion}

This paper evaluated 18 prompt-design variants across six open-weight,
instruction-tuned LLMs on the official EDOS sexism-detection benchmark
\cite{kirk2023semeval}, with paired significance testing, subtype-level
error analysis, order, few-shot-size, and demonstration-selection
sensitivity sweeps, and a head-to-head comparison against QLoRA, an
encoder, and a classical TF-IDF baseline. The clearest result in this
study cuts against prompting. Most of this study's prompt-design choices produce
large, statistically real effects (92 of 101 comparisons significant
under McNemar's test, 88 of 101 under the paired F1-diff bootstrap), and
phi3's largest single effect (+0.428 F1) is the biggest shift this study
puts through significance testing, though not the biggest shift in the
study overall: QLoRA fine-tuning moves the same model further still
(Sec.~\ref{sec:discussion:prompt-effects}). But
those effects operate entirely within a ceiling that a linear classifier
over word counts already clears for five of the six models, at a
fraction of the computational cost. Only qwen3 beats that classifier on
F1, and it does so at 4--16$\times$ the inference cost of every other
prompted model, while still losing to it on accuracy. Prompt design is a
real, measurable lever on this task, but the evidence here does not make
it a substitute for adapting a model, however simple, to the task's own
labeled data.

Two further findings deserve separate mention. Llama-3.1-8B's
interaction with chain-of-thought prompting is not the single-variant
curiosity it first appeared to be: every chain-of-thought variant this
model was evaluated on shows elevated generation failures, and the cause
splits cleanly depending on whether the prompt also asks for
aspect-by-aspect analysis, explicit safety refusal when it does not,
reasoning-token-budget exhaustion when it does
(Sec.~\ref{sec:results:llama-refusal}). Neither failure mode was visible
in the aggregate accuracy or F1 numbers alone; both surfaced only once
the raw generations were read directly. A model's headline score can
look ordinary while concealing a specific, systematic reason it is
wrong. Separately, prompt-\emph{component} order turned out to matter
far less than prior work on few-shot demonstration order would predict,
moving only one of six models' scores in either direction
(Sec.~\ref{sec:results:sensitivity}). In this study, which components a
prompt includes matters far more than the order they appear in.

For anyone building a sexism-detection or similar content-moderation
pipeline on an off-the-shelf instruction-tuned LLM, prompt engineering
is not wasted effort: it measurably changes what these models can do.
But it should be budgeted and evaluated as one option among several
rather than assumed to be the strongest one, and any pipeline that
treats a model's silence or refusal as an ordinary negative prediction
should monitor for that directly instead of trusting an aggregate score
to reveal it. The most
direct next steps are the ones this paper's limitations already name:
testing whether the same pattern holds for closed, commercial models,
and extending confidence-based scoring to chain-of-thought variants,
the one part of this study's own grid, including qwen3's best-performing
variant, that first-token AUC cannot reach at all.


\bibliographystyle{splncs04}
\bibliography{references}

\end{document}
